\documentclass[10pt]{article}
\usepackage[margin=20mm]{geometry}
\usepackage{amsmath,amssymb,mathtools,fontspec,unicode-math,xurl,hyperref}
\IfFontExistsTF{DejaVu Serif}{\setmainfont{DejaVu Serif}}{\setmainfont{Latin Modern Roman}}
\setmathfont{latinmodern-math.otf}
\setlength{\parindent}{0pt}\setlength{\parskip}{6pt}\sloppy
\title{Scientific initial review of M25-0018}
\author{}
\date{}
\begin{document}
\maketitle
Abstract: The acquired September 2024 manuscript asserts an effective equidistribution alternative for a unipotent flow on the space of unimodular three-dimensional lattices. An explicit rational lattice satisfies the stated two-sided compactness test with a strict factor-two margin, while every point of the tested orbit contains the same fixed vector. A compactly supported test function therefore contradicts the main theorem. The same example makes the initial deletion step remove the whole orbit. A separate small-window calculation disproves the parameter-recovery estimate in Lemma 11.4. A short-time separation claim also needs a restricted range. All 63 source pages and the complete matching TeX and bibliography were inspected. These conclusions concern arXiv version 6 only; no mathematical equivalence to the journal final has been established.\par
\section*{Identity, version, and scope}
Lei Yang, Effective version of Ratner’s equidistribution theorem for \(\operatorname{SL}(3,\mathbb R)\), associated journal DOI 10.4007/annals.2025.202.1.3 (https://doi.org/10.4007/annals.2025.202.1.3). The actual scientific input is arXiv:2208.02525v6 (https://arxiv.org/abs/2208.02525v6), dated 19 September 2024, 63 pages. All page references below are printed manuscript pages, equal to one-based PDF pages. Publisher metadata gives journal pages 189–264; that identity information does not establish equivalence of the acquired 63-page manuscript and the final publication.\par
The PDF `papers/0018/0018.pdf`, 554,813 bytes, has SHA256\par
> 6ad6199b4aaf27439d4ee3baa0f3b3bb86d8809aa969a6604469ecab4178ba0a.\par
The source manifest has SHA256\par
> 07f97f1737fca2620404f36ec4c0bca414c3389e2213d78e9039615e750db677.\par
All eight manifest-bound files were rehashed successfully. The matching source tree contains `main.tex` (3,360 lines) and `main.bbl` (201 lines). Both were read to the end. No separate computational code release or standalone supplement was acquired. Author code and author TeX were not executed.\par
\section*{Checked claims and outcome}
The initial review read Sections 1–12 and the complete bibliography. It traced the dependence of Theorems 1.2–1.3 on the initial nondivergence/deletion step, transverse dimension improvement, and the high-dimension equidistribution input. It inspected Corollaries 1.4 and 1.6, the Lie-algebra and coordinate calculations, all stated scale ranges, the heavy/unstructured/structured decompositions, and the final measure approximation.\par
The supported outcome is Major: the acquired main theorem is false as stated. Problem 1 is a complete counterexample, with both primal and exterior-square conditions checked for all relevant times. Problem 2 independently invalidates a core counting estimate; it does not by itself show that every subsequent counting conclusion is false. Problem 3 is a local short-time defect whose later application uses large times. Three other transitions remain Needs Evidence and are not counted as supported defects.\par
Pages | Claims inspected | Result/limit\par
1–5 | Definitions of \(X_\varpi\), \(a,a_0,U\), Theorems 1.2–1.3, line and curve applications | Complete counterexample to the main alternative; graph-coordinate and Taylor issues distinguished.\par
6–12 | Lie-algebra decomposition, BCH estimates, high-dimension input, nondivergence and separation | Basic brackets and weights checked; Problem 3. External mixing/nondivergence proofs not independently reproduced.\par
13–17 | Initial set \(\mathcal F\), closing/dimension bounds, iteration and finite-cover reduction | Problem 1 makes \(\mathcal F\) empty. Several explicit scope/sign corrections recorded separately.\par
18–30 | Flow boxes, projections, tube coordinates, intersections and hyperbolas | Coordinate formulas and their domains inspected; cross-section boundary control remains pending.\par
31–43 | Heavy, strange, unstructured and coarse-scale dimension estimates | Exponents and normalization dependencies traced; later normalization factors provide counterevidence to one earlier omission.\par
44–59 | Structured part, parameter recovery and curve counting | Problem 2; full downstream repair not verified.\par
59–63 | Final decomposition, iteration, theorem proof and references | Scale budget checked; exact versus weak measure equality remains pending.\par
\section*{Definitions and quantifiers}
Let \(G=\operatorname{SL}(3,\mathbb R)\), \(\Gamma=\operatorname{SL}(3,\mathbb Z)\), and \(X=G/\Gamma\). A point \(g\Gamma\) represents the lattice \(\Lambda=g\mathbb Z^3\). The manuscript defines \[u(r)=I+rE_{23},\qquad
 a(t)=\operatorname{diag}(1,e^{t/2},e^{-t/2}),\qquad
 a_0(\ell)=\operatorname{diag}(e^{\ell/3},e^{-\ell/6},e^{-\ell/6}).\] Its compact set is \[X_\varpi=\{g\Gamma:\|gv\|\geq\varpi
 \text{ for all }0\ne v\in\mathbb Z^3\text{ and all }0\ne v\in\wedge^2\mathbb Z^3\}.\] The exterior-square condition is essential: checking only lattice vectors would not test the theorem. For an oriented unimodular Euclidean three-dimensional lattice, the Hodge map identifies \(\wedge^2\Lambda\) isometrically with the dual lattice \(\Lambda^*=g^{-T}\mathbb Z^3\).\par
Equation reference: eq:main1, eq:main2. Theorem 1.3 on page 3 asserts that fixed positive \(\eta,t_0\) work for every \(t\geq t_0\) and every \(x_0\in X\), with one of the following alternatives: \[\begin{aligned}
 \left|\int_{-1/2}^{1/2}f(a(t)u(r)x_0)\,dr-\int_Xf\,d\mu_G\right|
 &\ll e^{-\eta t}\|f\|_S\quad\hbox{for every }f\in C_c^\infty(X),\\
 a_0(\ell)x_0&\notin X_{e^{-|\ell|/3}}
 \quad\hbox{for some }0.99t\leq|\ell|\leq1.01t.\end{aligned}\] The implicit constant is absolute under the convention on page 2. The Sobolev norm is fixed. Even allowing a finite constant depending on the single fixed lattice below would not repair the contradiction as \(t\to\infty\).\par
The notation \([L]\) means \([-L/2,L/2]\). Section 5 fixes \[\delta_1=\delta_2=10^{-5},\ d_0=10^{-6},\ C_1=10^7,\quad
 0<\theta<10^{-10},\quad \varepsilon_2=10^{-26}\theta,\quad N_1=10^5/\varepsilon_2,\] \[\varepsilon=e^{-10^5N_1},\quad\beta=e^{-\varepsilon^3t},\quad
 \delta_3=10^{-5}\theta\varepsilon,\quad \varepsilon_1=2\varepsilon_2.\] The iteration uses \(s_i=2^{-i}\delta_1t/2\) and \(d_i=d_0+i\varepsilon_2\). These constants, however extreme, are positive and fixed before \(t\) tends to infinity. No floating-point underflow is used to replace them by zero.\par
\section*{Supported substantive problems}
Problem 1 (Major: an explicit lattice contradicts the main alternative).  Location and suspect implication. Theorems 1.2–1.3, page 3; Definition 6.1 and equation (6.7), page 13; final proof, pages 61–62. The two-sided \(a_0\) test in [eq:main2] (\#eq:main2) does not exclude an \(AU\) orbit lying in a proper fixed-vector locus. Consequently the claimed small deletion and the subsequent dimension argument cannot establish [eq:main1] (\#eq:main1).\par
Parameters and every compactness condition. Set \[g=\begin{pmatrix}2&1/2&0\\0&1&0\\0&0&1/2\end{pmatrix},
 \qquad x_0=g\Gamma.\] Then \(\det g=1\), and \[\begin{aligned}
 \Lambda&=\{(2p+q/2,q,r/2):p,q,r\in\mathbb Z\},\\
 \Lambda^*&=\{(p/2,q-p/4,2r):p,q,r\in\mathbb Z\}.\end{aligned}\] The inverse transpose is \[g^{-T}=\begin{pmatrix}1/2&0&0\\-1/4&1&0\\0&0&2\end{pmatrix}.\] These formulas are exact, including the denominator four in the dual.\par
For \(L\geq12\log2\), put \(z=e^{L/6}\geq4\). The compactness threshold is \(e^{-L/3}=z^{-2}\). After dividing all coordinates by that threshold, the vectors in the four lattices to be checked are: \[\begin{array}{c|c}
 a_0(L)\Lambda &(z^4(2p+q/2),\,zq,\,zr/2)\\
 (a_0(L)\Lambda)^* &(p/2,\,z^3(q-p/4),\,2z^3r)\\
 a_0(-L)\Lambda &(2p+q/2,\,z^3q,\,z^3r/2)\\
 (a_0(-L)\Lambda)^* &(z^4p/2,\,z(q-p/4),\,2zr).
\end{array}\] For every nonzero integer triple, the maximum absolute coordinate is at least two:\par
1. In the first row, if \(2p+q/2\ne0\), its magnitude is at least \(1/2\), so the first coordinate is at least \(z^4/2\geq128\). If it is zero, at least one of \(q,r\) is nonzero, and the last two coordinates have maximum at least \(z/2\geq2\).\par
2. In the second row, if \((q-p/4,2r)\ne(0,0)\), the maximum of its coordinates is at least \(1/4\), so the corresponding maximum is at least \(z^3/4\geq16\). Otherwise \(r=0,p=4q\), and nonzeroness gives \(|p/2|=2|q|\geq2\).\par
3. In the third row, if \((q,r)\ne(0,0)\), the maximum of the last two coordinates is at least \(z^3/2\geq32\). Otherwise \(|2p|\geq2\).\par
4. In the fourth row, \(p\ne0\) gives a first coordinate at least \(z^4/2\). If \(p=0\), then \((q,r)\ne(0,0)\) and the maximum of the last coordinates is at least \(z\geq4\).\par
Equation reference: eq:strict. Thus, in both the usual Euclidean and maximum-coordinate norms, \[a_0(L)x_0,\ a_0(-L)x_0\in X_{2e^{-L/3}}\subset X_{e^{-L/3}}
 \qquad(L\geq12\log2).\] The exterior-square checks are the second and fourth rows, not omitted assumptions. Therefore [eq:main2] (\#eq:main2) is false for every allowed \(\ell\) once \(t\geq12\log2/0.99\). This is a strict factor-two example, not an issue about equality at the boundary of \(X_\varpi\).\par
Failure of the equidistribution alternative. The vector \(v=2e_1\) lies in \(\Lambda\), and \(a(t)u(r)v=v\) for all real \(r,t\). Hence every lattice \(a(t)u(r)x_0\) lies in \[S_v=\{\Lambda'\in X:v\in\Lambda'\}.\] This set is closed. Indeed, in a local chart choose representatives \(g_n\to g\) for a convergent sequence of lattices in \(S_v\). If \(g_nk_n=v\) with \(k_n\in\mathbb Z^3\), then \(k_n=g_n^{-1}v\) stays bounded. A subsequence is constant, so \(v\in g\mathbb Z^3\).\par
The set is proper: the unimodular diagonal lattice represented by \(\operatorname{diag}(\sqrt2,1,1/\sqrt2)\) does not contain \(2e_1\), since \(2/\sqrt2\) is not an integer. Choose a nonzero nonnegative smooth compactly supported function \(f\) in a relatively compact open subset of \(X\setminus S_v\). Such a function exists in a manifold chart. Haar measure has positive mass on every nonempty open subset, so \[I:=\int_X f\,d\mu_G>0,\qquad
 \int_{-1/2}^{1/2}f(a(t)u(r)x_0)\,dr=0
 \quad\hbox{for all }t.\] The left side of [eq:main1] (\#eq:main1) is the fixed positive number \(I\), whereas its proposed upper bound tends to zero. Taking \(t\geq\max(t_0,12\log2/0.99)\) and then \(t\) sufficiently large contradicts both alternatives. Theorem 1.2 is contradicted by the same calculation with its starting point \(a(t)x_0\), exactly as in the equivalence given on page 3.\par
The first failed proof step is also explicit. Write \(s=(1-\delta_1)t\) and \(\mathcal E=a(s)u([1])x_0\). Every \(y\in\mathcal E\) contains \(2e_1\). Definition 6.1 permits \(\ell=0\) and \(\ell'=-2\delta_1t\). Then \(a_0(\ell')a(-\ell)y\) contains \[2e^{-2\delta_1t/3}e_1.\] For \[t>\frac{\log2}{2\delta_1/3-\varepsilon^3/4},\] whose denominator is positive by the listed constants, its length is less than \(\beta^{1/4}=e^{-\varepsilon^3t/4}\). Every point is therefore removed: \(\mathcal F=\varnothing\). The removed proportion is one, while the bound claimed before (6.7) is \(O(s^2\beta^{\alpha_2})\to0\). The normalized measure on \(\mathcal F\) is not defined. This identifies an actual failure in the proof, beyond the abstract theorem counterexample.\par
Counterevidence, alternative readings, and repair limits. The familiar lattice \(\mathbb Z^3\) would already expose a boundary-equality concern, but it is unnecessary: [eq:strict] (\#eq:strict) has a margin. Including the dual lattice does not rescue the statement. Replacing a norm by another fixed norm does not rescue a uniform theorem either. For a positive integer \(N\), use \[g_N=\begin{pmatrix}N&1/N&0\\0&1&0\\0&0&1/N\end{pmatrix},
 \qquad
 g_N^{-T}=\begin{pmatrix}1/N&0&0\\-1/N^2&1&0\\0&0&N\end{pmatrix}.\] Both the primal and dual intersections with the first coordinate axis have primitive generator \(Ne_1\). All nonaxis normalized vectors eventually grow under the same coordinate cases. Choosing \(N\) large handles any fixed equivalent-norm constant. This also shows that merely inserting a fixed multiplicative constant in the threshold cannot uniformly repair the theorem over all starting lattices.\par
A slower decay threshold \(e^{-(1/3-\kappa)|\ell|}\) with fixed \(\kappa>0\) would eventually detect the fixed-axis examples, but this observation is not a proof of a corrected theorem. A valid replacement must establish the quantitative implication needed for Definition 6.1, with all primal and dual obstruction heights controlled. No such global repair is claimed here.\par
Impact and evidence. This is material damage to the central theoretical guarantee, hence Major under 0.7.1. It does not refute Ratner’s qualitative theorem: the obstruction is precisely that the orbit remains in a proper locus. The proof is analytic and covers infinitely many integer vectors and all sufficiently large times. The reviewer-owned script separately verifies determinant, dual, Hodge matrix, fixed vector, and 58,944 illustrative finite-box inequalities; finite enumeration is not substituted for the coordinate proof.\par
Problem 2 (Major: small-window proximity does not give (11.2)).  Location and dependency. Lemma 11.4(1), equation (11.2), pages 44–46, asserts parameter recovery from proximity of two hyperbolas in \(\mathfrak C=[1.2\beta e^{-s'}]^2\). It is explicitly applied on pages 49 and 51 in Lemmas 11.5 and 11.7, which enter Proposition 11.1 and hence the central dimension-improvement Proposition 9.1. The proof replaces the available \(\beta e^{-s'}\) window by an \(e^{-s'}\) scale when extracting coefficients.\par
Full admissible construction. Write \(S=s'>0\), \(E=e^S\), and \(h=C_1\varepsilon_1\). Choose the manuscript parameters at sufficiently large \(t\) so that \[0<\beta\leq10^{-2},\quad E^2\geq4\beta^{-2},\quad
 e^{-10hS}\leq\beta/4.\] Take \[\ell=0,\quad k=k'=1,\quad A=1,\quad
 w=\beta E/4,\quad w'=w+\beta^{-1}E^{-1}.\] Then \(|k|\in[e^\ell,e^{\ell+1}]\), \(0\leq\ell\leq S\), and \[e^{S-10hS}\leq w<w'\leq\tfrac12\beta E<\beta E.\] Also \(0<A<e^{(2-h)S}\) because \(0<h<2\). Thus every parameter restriction in part (1) is met.\par
These are compatible with the actual hierarchy, not merely freely chosen surrogate parameters. Indeed \(S=ct\) at an iteration scale, with \(c\geq10^{-5}2^{-N_1-1}\) for \(i\leq N_1\). Since \(\varepsilon^3=e^{-3\cdot10^5N_1}\), \[c>\varepsilon^3,\qquad 10C_1\varepsilon_1c>\varepsilon^3.\] For the second inequality, its left side at the displayed lower bound for \(c\) is \((10^8/N_1)2^{-N_1}\), much larger than \(e^{-300000N_1}\). These inequalities can be checked by logarithms, using \(N_1>10^{41}\) and \(\log N_1<N_1\). They ensure all three conditions above once \(t\) is large. In particular no forbidden interchange of constants and time is used.\par
Let \[C=\{(X,Y)\in\mathfrak C:X=-Y/(1+wY)\},\qquad
 C'=\{(X,Y)\in\mathfrak C:X=-Y/(1+w'Y)\}.\] For \(|Y|\leq0.6\beta/E\), both denominators are at least \(1-0.3\beta^2\). At the same \(Y\) the horizontal separation is exactly \[|X'-X|
 =\frac{\beta^{-1}E^{-1}Y^2}{|(1+wY)(1+w'Y)|}
 \leq\frac{0.36\beta}{(1-0.3\beta^2)^2}E^{-3}
 <E^{-3}.\] For \(Y\geq0\), the corresponding point of \(C\) lies in the box because \(|X|\leq Y\). For \(Y<0\), \(w'>w\) gives \(0<X<X'\), so whenever \((X',Y)\) lies in the box, \((X,Y)\) does too. Thus the entire truncated \(C'\) lies in the asserted \(AE^{-3}\) Euclidean neighborhood of \(C\), including near the box boundary.\par
The conclusion (11.2), however, requires \(|w-w'|\leq AE^{-1+0}=E^{-1}\), whereas \[|w-w'|=\beta^{-1}E^{-1}.\] The ratio \(\beta^{-1}=e^{\varepsilon^3t}\) is unbounded. An absolute implicit constant cannot cover the discrepancy. This is not a boundary, sign, zero-denominator, or \(k'=0\) example.\par
Counterevidence and a precisely limited correction. The subsequent arguments allow large exponential slack, and the hierarchy has \(\log\beta^{-1}\ll\varepsilon_1s'\) throughout the required finite iteration. It is therefore plausible that explicitly retaining a fixed power of \(\beta^{-1}\) could be absorbed later. That possibility does not validate the displayed implication or establish the repaired counting argument.\par
In a restricted graph regime where \(k,k'\) are comparable to \(e^\ell\), denominators stay comparable to them, and the horizontal (same-\(Y\)) discrepancy is at most \(\rho\) on two points \(Y=\pm b\) with \(b\asymp\beta e^{-S}\), the exact identity gives \[|\Delta k\pm b\Delta w|\leq C\rho e^{2\ell}/b.\] Adding and subtracting yields the defensible local bounds \[|\Delta k|\leq C\rho e^{2\ell}/b,\qquad
 |\Delta w|\leq C\rho e^{2\ell}/b^2.\] With \(\rho=Ae^{-3S}\) these contain respectively \(\beta^{-1}\) and \(\beta^{-2}\). This is a conditional elementary correction, not a replacement of all parts of Lemma 11.4. A full repair must also justify denominator comparability, conversion from Euclidean neighborhood to a same-\(Y\) horizontal discrepancy bound, truncated-curve endpoints, and propagation through Lemmas 11.5–11.7.\par
Impact and evidence. The invalid coefficient estimate is a demonstrated derivation error in the core dimension argument and is Major even if a corrected exponent budget ultimately works. No separate counterexample to every downstream counting bound is claimed. The exact rational separation identity and the uniform small-window upper bound are verified in the reviewer-owned script; the all-time admissibility argument is analytic.\par
Problem 3 (Minor: Lemma 4.2 includes a short time excluded by its proof).  Location and claim. Lemma 4.2 on page 12 is stated for every \(\ell\geq1\). Under its injectivity and coordinate hypotheses it concludes that the negative-weight coefficient pair has norm at least \(\beta e^{-\ell/2}\). The proof first uses \(\ell\) sufficiently large to exclude the identity lattice element; that step is false at the stated lower endpoint.\par
Complete check in the standard coordinate norm. Choose \(0<\beta<10^{-10}\), \(\ell=1\), \(x=a(1)\Gamma\), \[r_1=r_3=0,\qquad r_2=\beta/4,\qquad
 \mathbf v=e^{1/2}\beta E_{21}/4.\] Here \(\mathbf v_1=E_{21}\) as defined on page 6. The parameters lie in \([\beta]\) and their Euclidean triple norm is \(\beta/4\). The coefficient norm of \(\mathbf v\) is \(e^{1/2}\beta/4<\beta\); this also holds for the Frobenius norm and twice that norm. Moreover \(a(-1)x=\Gamma\in X_{\beta^{1/4}}\). The equality of the two exponential actions in the hypothesis is an identity of group elements: \[\exp(r_1e^\ell\mathbf u+r_2e^{\ell/2}\mathbf v_1+
                   r_3e^{\ell/2}\mathbf w_1)=\exp(\mathbf v).\] The \(\mathbf u^*\), \(\mathbf v_2\), and \(\mathbf w_2\) coefficients of \(\mathbf v\) are all zero. Thus \(|u^*|<\beta e^{-1}\), while the claimed conclusion would read \(0\geq\beta e^{-1/2}>0\). In the proof the lattice element is exactly the identity.\par
Counterevidence, repair and impact. The source does not specify an exceptional rescaling of the Lie-algebra norm that would exclude these usual conventions. Any fixed-norm formulation should state the required lower time depending on that convention. The later application in Section 6 uses time proportional to a large parameter, so this example does not establish failure of that application. Restricting the lemma to sufficiently large \(\ell\) removes this identity-element example; a complete proof of the nonidentity case and its constants still has to be supplied. This is a local Minor range defect, independent of the much more serious failure in Problem 1.\par
\section*{Needs Evidence: unresolved transitions}
These items are not additional supported problems. The checks below identify what is missing without claiming counterexamples to conclusions whose full hypotheses have not been tested.\par
\subsection*{NE1: the curve linearization needs a smaller scale or a different argument}
Corollary 1.6, page 5, divides a smooth nondegenerate graph into pieces of length \(e^{-t/2}\) and claims exponentially small error after applying \(a_1(t)\), where \(a_1(t)=\operatorname{diag}(e^{t/3},e^{t/3},e^{-2t/3})\). For the admissible graph \(\psi(r)=(r,r^2)\), the first and second derivatives span \(\mathbb R^2\) at every point. At an increment \(\Delta=q e^{-t/2}\) the quadratic Taylor remainder is \(\Delta^2E_{13}\) before conjugation. Since \(\operatorname{Ad}(a_1(t))E_{13}=e^tE_{13}\), the left group-coordinate remainder is exactly \[e^t\Delta^2E_{13}=q^2E_{13}.\] For fixed nonzero \(q\in[-1/2,1/2]\) it does not tend to zero.\par
This disproves the direct exponentially small group-Taylor estimate, but it has not been turned into a uniform lower bound for quotient distance or for the tested orbit integrals: cusp identifications and the relevant norms must be controlled. It is therefore recorded as Needs Evidence concerning the corollary’s derivation, not as a counterexample to the corollary. Using pieces of length \(e^{-(1/2+c)t}\) gives the precise group-coordinate bound \(O(e^{-2ct})\) for bounded second derivative, but then the expansion time, diagonal factors and nondivergence argument all change. No completed repair of the corollary is claimed.\par
Two additional proof obligations belong to this same application: a general nondegenerate curve need not globally be the displayed graph without charts and Jacobian weights; and controlling each exceptional set for a fixed real \(\ell\) does not itself control their uncountable union. An explicit time mesh and a compactness-threshold margin can potentially address the latter. The manuscript’s stronger \(e^{-t/10}\) threshold supplies useful slack, which is counterevidence against treating the union step as an immediate false theorem.\par
\subsection*{NE2: exact positive decomposition versus weak approximation}
Proposition 12.1 and (12.1), pages 59–61, state an exact positive measure decomposition after discarding a small positive part. The construction then projects orbit pieces and proves only a \(C^1\)-dual estimate (12.5), of size \(O(\beta)\). A small weak error need not be a positive discarded measure. Also, normalizing restrictions to an overlapping cover does not by itself yield a probability decomposition with the displayed weights; disjoint cells or a partition of unity and the induced weights must be specified.\par
The final theorem proof on page 61 actually uses an \(O(\beta)\) remainder. That is significant counterevidence: a weak approximation may suffice for the intended application. The remaining obligation is a precise decomposition in the correct topology, with masses summing appropriately, overlap controlled, and the projected transverse measures satisfying the needed dimension bound. No counterexample satisfying all of Proposition 12.1’s dimension assumptions was constructed here, so this remains Needs Evidence.\par
\subsection*{NE3: cross-section boundary control}
Lemma 8.5, pages 20 and 23, claims uniformly long intersections with every cross-section for every point of the full playground. Its proof controls the longitudinal coordinate \(f_y(r)\), but membership in the cross-section also imposes transverse-coordinate bounds. The displayed formulas (8.8)–(8.10) show that these transverse coordinates move. For example, in the subcase \(w_2=0\), the positive coordinate moves as \(v_1(r)=v_1+rv_2\); a point close to a box face can leave that face while the longitudinal coordinate is brought to zero.\par
The relevant global injectivity, boundary trimming and full-playground hypotheses were not all established for a single proposed counterexample, so no failed lemma is asserted on that basis. The needed evidence is a nested/fattened-box formulation or a quantified boundary deletion which proves the uniform intersection lengths actually used in Lemmas 8.6 and 8.13. Merely solving the longitudinal equation does not supply that evidence.\par
\section*{Scope clarifications and counterevidence retained}
\subsection*{The sign in Proposition 4.1 is not counted as a separate substantive finding}
On page 11 the hypothesis literally says both \(a_0(\pm s)x\) are outside \(X_{e^{-s/3}}\), whereas the contradiction proof needs them inside. The literal statement is false: for \(s=1\) and \(g_\delta=\operatorname{diag}(\delta,1,\delta^{-1})\), both tests are outside for sufficiently small \(\delta\), and every \(a(1)u(r)g_\delta\Gamma\) contains \(\delta e_1\). Taking \(\beta=\sqrt\delta\) makes the exceptional proportion one and the claimed \(O(\beta^\alpha)\) bound tend to zero. The proof and later intended use identify the sign correction. As an apparent negation transcription error it is recorded here, not used to inflate the substantive count. Correcting that sign does not repair Problem 1 or (6.7).\par
\subsection*{Normalization in Lemma 8.6}
Page 21 defines \(\widetilde\mu=\mu|_{\Pi}/\mu(\Pi)\) but bounds a normalized mass by the corresponding unnormalized mass. For a measurable subset \(A\subset\Pi\) of positive mass and \(p=\mu(\Pi)<1\), the exact relation is \(\widetilde\mu(A)=\mu(A)/p>\mu(A)\). The available full-playground condition gives \(p\geq\beta^{15}\) and therefore a factor \(\beta^{-15}\), not factor one. The manuscript itself includes this factor in Lemma 9.4 on page 32 and again in the application on page 39. These later formulas are the relevant counterevidence and indicate the intended correction. No additional core failure is inferred solely from the earlier missing factor, and the exact unchanged statement of Lemma 8.6 is not declared proved.\par
\subsection*{A finite cover must cover a compact relevant region}
The first line of the proof of Proposition 7.2, page 17, covers the whole noncompact space \(X\) by finitely many translates of a bounded local box. A finite union of their compact closures cannot cover \(X\). The trajectory being studied starts from a compact parameter interval, and the argument also retains thick points; a relevant compact region can replace \(X\). The required cardinality and overlap bounds still need an actual covering argument on that region. This is a scope clarification, not a claim that compact truncation is impossible.\par
\subsection*{Greedy counting can use first ownership}
On page 57 the greedy construction is described as giving each selected curve many cells not covered by any other selected curve. The construction directly ensures only that its newly added cells were not covered by earlier selections. For a lower bound on the union, assign each cell to the first selected curve that covers it. These assigned new-cell sets are disjoint, so their cardinalities add. This repairs that local interpretation of the greedy union count and does not address the parameter estimate in Problem 2.\par
\subsection*{Correct scale and algebra checks}
The reviewer checked \([\mathbf u,\mathbf v_2]=\mathbf v_1\), \([\mathbf u,\mathbf w_2]=-\mathbf w_1\), and \([\mathbf v_1,\mathbf w_1]=\mathbf u\) directly. The elapsed-time identity \[(1-\delta_1)t+\sum_{j=0}^{i-1}s_j=t-2s_i\] is correct. There is also adequate formal room in the stated structured-dimension range: \[10^{11}C_1\varepsilon_1=2\cdot10^{-8}\theta<\theta.\] Thus the restriction in Proposition 11.1 is compatible with stopping at dimension \(4-\theta\). These checks do not validate the missing initial dimension input. They prevent an incorrect additional allegation about the iteration budget.\par
\section*{Dependence, version limits, and repair priorities}
Problem 1 defeats the main theorem before any of the detailed Kakeya estimates can restore it. The incorrect implication is from the displayed \(a_0\) compactness test to a predominantly retained set \(\mathcal F\) with transverse dimension. The complete deletion calculation locates this failure at the entrance to the proof. An initial theorem repair must precede reassessing the iterative dimension estimates and the two advertised applications.\par
Problem 2 is independently relevant to the proof method. A new theorem with a corrected obstruction condition would still require a valid coefficient-recovery lemma, its admissible ranges, and a checked downstream exponent budget. Large spare constants are not a proof that every use survives. Problem 3 has a limited scope: its small-time counterexample is not asserted to affect the large-time closing application.\par
A bounded primary-source search checked the arXiv metadata and publisher identity page and searched the exact title/identifier with correction and erratum terms. The metadata returned version 6 dated 19 September 2024; no relevant correction was located in those returned results. This is not evidence that no correction exists. No final journal PDF, later private revision, or author correspondence was inspected. All findings must be compared with the final publication before being described as final-publication defects.\par
There are no empirical experiments to reproduce in the acquired materials. Required next work is mathematical: formulate the correct obstruction, prove its quantitative deletion consequence, repair the local curve estimates, and supply the indicated boundary/decomposition arguments. Running author code is neither necessary nor authorized.\par
Item | Judgment | Repair verification | Adoption/remaining work\par
P1 | Major; main-theorem counterexample | No corrected theorem proved; fixed-prefactor repair excluded by \(g_N\) family | Author adoption unrequested; new obstruction theorem needed.\par
P2 | Major; core derivation invalid | Conditional two-point coefficient estimate proved; global counting repair incomplete | Recheck all uses and exponent margins.\par
P3 | Minor; short-time range | Identity example completely verified; sufficiently large-time restriction removes this example only | Full revised lemma proof pending.\par
NE1 | Needs Evidence | Taylor residual computed; smaller group-Taylor scale identified | Quotient/integral argument and corollary proof pending.\par
NE2 | Needs Evidence | Exact and weak decompositions distinguished | Positive weights/cover/weak-error statement pending.\par
NE3 | Needs Evidence | Moving transverse coordinates identified | Admissible global example or boundary-control proof pending.\par
Scope | Clarifications | Sign, normalization, finite-cover scope and greedy ownership separated | Not additional Major findings.\par
\section*{Consulted sources}
The scientific source consulted in full is Yang’s version-bound manuscript and its complete matching TeX/bibliography. The metadata pages below were used for identity/version information. The paper’s citations to quantitative nondivergence, mixing and effective closing are dependencies inspected as invoked; their complete external proofs were not independently reproduced in this initial review. No statement is certified merely because it is cited.\par
9 L. Yang, Effective version of Ratner’s equidistribution theorem for \(\operatorname{SL}(3,\mathbb R)\), arXiv:2208.02525v6, 19 September 2024, 63 pp. https://arxiv.org/abs/2208.02525v6. Manifest-bound PDF and source archive are the full reviewed scientific materials. arXiv metadata and submission history for 2208.02525, https://arxiv.org/abs/2208.02525. Version/date check only; retrieved through the available search tool. Annals of Mathematics, identity page for DOI 10.4007/annals.2025.202.1.3, https://annals.math.princeton.edu/2025/202-1/p03. Metadata only; journal-final mathematical equivalence not established.\par
\section*{Actual independent checks and evidence}
The reviewer wrote and executed only `verification/0018/evidence/independent\_checks.py`. It uses the existing Python 3.13.5 and SymPy 1.13.3 installation, without network access, package installation, or author-code execution. The actual command was PowerShell invocation of\par
> `[local archive path omitted] `-X utf8 -B`
> `[staging]\textbackslash{}verification/0018/evidence/independent\_checks.py`.\par
Process 42108 exited with code zero and reported 12 passing check groups and 58,944 finite lattice evaluations. The exact script and JSON results are retained. Check groups comprise determinant/dual, Hodge identification, fixed vector, finite lattice illustrations, the \(g_N\) family, optional period \(u(8)\), exact hyperbola separation, its small-window bound, the curve Taylor residual, the short-time separation example, three Lie brackets, and the structured-range slack.\par
The optional identity \(g^{-1}u(8)g\in\operatorname{SL}(3,\mathbb Z)\) checks a period of the starting unipotent orbit. The main counterexample does not depend on a closed-orbit theorem or on this optional calculation. The finite enumeration ranges over \(-8\leq p,q,r\leq8\), excluding zero, for \(z=4,8,16\) and all four scaled lattices. Its finite nature is explicitly recorded; the report’s case split is the universal proof.\par
Static inspection of `main.tex` and `main.bbl` records their hashes, line counts, section declarations, citation keys and file/command references. It is a read-only text inspection. It is not a TeX build, executable audit of unavailable code, proof-assistant verification, or validation of every cited paper.\par
\section*{Source reading, page inspection, and compilation failure}
All 3,360 lines of `main.tex` were directly read in ten contiguous ranges: 0–349, 350–699, 700–1049, 1050–1399, 1400–1749, 1750–2099, 2100–2449, 2450–2799, 2800–3149, 3150–3359. The last read returned zero remaining lines. All 201 bibliography lines were read. Extracted page text was produced to support location checks; this report does not claim that the whole extracted-text file was separately reread in addition to the source.\par
All 63 original PDF pages were rendered with Poppler and actually viewed. Page images checked the theorem signs, matrices, compactness/exterior-square conditions, scale exponents, hyperbola figure, cross-section formulas and reference list. The renderer reported a warning that no display font for Symbol was found. The inspected mathematical content was visible in the images and checked against the matching TeX; the warning remains in the actual render log and is not suppressed.\par
The Windows capability check found no `pdflatex`, `xelatex`, `lualatex`, or `tectonic` on PATH or in the checked standard locations. No installation was performed. Static environment-balance checks are not a substitute for compilation. These missing delivery steps are recorded in `compilation\_status.json` and `delivery\_state.json`.\par
\section*{Supporting checks}
The matching `main.tex` SHA256 is\par
> dba77e6abb6f621de3f32d0a26b9b6dd553196436b9691b1ea10ab6266fe06c0;\par
the `main.bbl` SHA256 is\par
> f95d4828ee078fca3cb51d811502e8639f28bd36941ce0bf71e26f813465d63b.\par
That observation belongs to the dispatcher; this assistant did not personally observe the dispatcher’s browser tab. The two names and their sources are retained separately, without an equivalence claim or independent backend-route verification. No model switch was performed, and no independent evidence of a non-6-Pro backend route was observed.\par
All scientific reading, calculations, source inspection, page viewing and report writing described here were performed by this assistant. No scientific work was delegated. Human scientific approval, independent acceptance, journal endorsement and author adoption are not claimed. Only the assigned staging path was used for writes. No final archive, other-paper output, Library/cloud upload, author contact, heartbeat, package installation or author-code execution was performed.\par
\section*{Compilation continuation (26 October 2026)}
Scientific conclusions, pending questions, source-version limits and the absence of independent acceptance remain unchanged.\par
\end{document}
